\section{Exercises}

\subsection{Poisson Distribution}

\[
\text{PMF:} f(x \mid \theta)= \sum_{k=1}^K \pi_k * (\frac{\lambda_k^x}{x!} e^{-\lambda_k})
\]

\[
L(\theta) = \prod_{i=1}^n f(x_i \mid \theta) = \prod_{i=1}^{n} \sum_{k=1}^K \pi_k (\frac{\lambda_k^{x_i}}{x_i!} e^{-\lambda_k})
\]

\[
\ell(\theta) = \log L(\theta) = \sum_{i=1}^{n} \log(\sum_{k=1}^K \pi_k (\frac{\lambda_k^{x_i}}{x_i!} e^{-\lambda_k}))
\]

\[
L_c(\theta \mid X, Z) = \prod_{i=1}^n \prod_{k=1}^K [\pi_k f(x_i \mid \theta_k)]^{\tilde{Z}_{ik}}
\]

\[
= \prod_{i=1}^n \prod_{k=1}^K [\pi_k (\frac{\lambda_k^{x_i}}{x_i!} e^{-\lambda_k})]^{\tilde{Z}_{ik}}
\]
%log:
\[
\ell_c(\theta \mid X, Z) = \sum_{i=1}^n \sum_{k=1}^K \tilde{Z}_{ik} [\log \pi_k + \log f(x_i \mid \theta_k)]
\]

\[
= \sum_{i=1}^n \sum_{k=1}^K \tilde{Z}_{ik} [\log \pi_k + \log (\frac{\lambda_k^{x_i}}{x_i!} e^{-\lambda_k})]
\]

\[
= \sum_{i=1}^n \sum_{k=1}^K \tilde{Z}_{ik} [\log \pi_k + \log (\frac{\lambda_k^{x_i}}{x_i!}) + \log (e^{-\lambda_k}))]
\]

\[
= \sum_{i=1}^n \sum_{k=1}^K \tilde{Z}_{ik} [\log \pi_k + \log (\lambda_k^{x_i}) - \log (x_i!) + \log (e^{-\lambda_k}))]
\]

\[
= \sum_{i=1}^n \sum_{k=1}^K \tilde{Z}_{ik} [\log \pi_k + x_i\log \lambda_k - \sum_{j=1}^{x_i} \log j - \lambda_k]
\]

\subsection{Estimation of the Mixing Proportions $\pi_k$}

The mixing proportions satisfy the constraint
\[
\sum_{k=1}^K\pi_k=1.
\]

Equivalently,
\[
\sum_{k=1}^K\pi_k-1=0.
\]

We can write
\[
\sum_{k=1}^K\pi_k-1=\psi(x).
\]

To find the extremum of a function $\phi(x)$ under the constraint
\[
\psi(x)=0,
\]
we introduce a Lagrange multiplier $\nu$ and optimize
\[
F(x,\nu) = \phi(x)+\nu\psi(x).
\]

In our case, we therefore define the Lagrangian
\[
F(\theta,\nu) = \ell_c(\theta \mid X, Z) + \nu \left(\sum_{k=1}^K\pi_k -1 \right).
\]

We use the conditions
\[
\frac{\partial F}{\partial \pi_k}=0,
\qquad
\frac{\partial F}{\partial\nu}=0.
\]

When differentiating with respect to $\pi_k$, all terms that do not
depend on $\pi_k$ can be ignored. Therefore,
\[
F = \sum_{k=1}^K \left(\sum_{i=1}^n\tilde{Z}_{ik} \right)\ln \pi_k + \nu \left(\sum_{k=1}^K\pi_k -1 \right) + \text{terms independent of the }\pi_k.
\]

Define
\[
n_k = \sum_{i=1}^n\tilde{Z}_{ik},
\]
where $n_k$ is the number of observations belonging to group $k$.

\[
F = \sum_{k=1}^K n_k \log \pi_k + \nu \left(\sum_{k=1}^K\pi_k \right) -\nu + \text{constant}.
\]

Thus, the part of the Lagrangian that depends on the mixing
proportions becomes
\[
F = \sum_{k=1}^K n_k \log \pi_k + \nu \left(\sum_{k=1}^K\pi_k \right) + \text{constant}.
\]

\paragraph{Step 1: Differentiate with respect to $\pi_k$.}

For a fixed group $k$,
\[
F = \sum_{j=1}^K n_j \log \pi_j + \nu \left(\sum_{j=1}^K\pi_j \right) + \text{constant}.
\]

\[
\frac{\partial F}{\partial \pi_k} = n_k * \frac{1}{\pi_k} + \nu = \frac{n_k}{\pi_k} + \nu.
\]

At the optimum,
\[
\frac{n_k}{\pi_k}+\nu=0.
\]

Therefore,
\[
\pi_k=\frac{n_k}{-\nu}.
\]

\paragraph{Step 2: Differentiate with respect to $\nu$.}

The derivative with respect to the Lagrange multiplier gives back the constraint:
\[
\frac{\partial F}{\partial\nu} = \sum_{k=1}^K\pi_k - 1= 0.
\]

Therefore,
\[
\sum_{k=1}^K\pi_k=1.
\]

Substituting
\[
\pi_k=\frac{n_k}{-\nu}
\]
into the constraint gives
\[
\sum_{k=1}^K\frac{n_k}{-\nu}=1.
\]
Since $-\nu$ is independent of $k$,
\[
\frac{1}{-\nu}\sum_{k=1}^{K}n_k = 1
\]

And since every observation belongs to exactly one group,
\[
\sum_{k=1}^Kn_k=n.
\]

Therefore,
\[
\frac{n}{-\nu}=1
\qquad\Longrightarrow\qquad
\nu=-n.
\]

Hence,
\[
\boxed{
\hat{\pi}_k=\frac{n_k}{n}
}.
\]

Since
\[
n_k=\sum_{i=1}^n\tilde{Z}_{ik},
\]
we can equivalently write
\[
\boxed{
\hat{\pi}_k=\frac{1}{n}\sum_{i=1}^n\tilde{Z}_{ik}
}.
\]

Thus, the estimated mixing proportion of group $k$ is simply the
\hl{proportion of observations belonging to group $k$}.

\subsection{Estimation of $\lambda_k$}

Unlike the mixing proportions, $\lambda_k$ is not subject to a constraint. Therefore, we can directly differentiate the completed log-likelihood with respect to $\lambda_k$.

Starting from
\[
\ell_c(\theta \mid X, Z) = \sum_{i=1}^n \sum_{k=1}^K \tilde{Z}_{ik} [\log \pi_k + x_i\log \lambda_k - \sum_{j=1}^{x_i} \log j - \lambda_k]
\]

we only keep the terms that depend on $\lambda$
\[
\ell_c(\theta \mid X, Z) = \sum_{i=1}^n \sum_{k=1}^K \tilde{Z}_{ik} [x_i\log \lambda_k - \lambda_k] + \text{constant}
\]

\paragraph{Step 1: Differentiate with respect to $\lambda_k$.}

For a fixed group $k$, we have
\[
\ell_c(\theta \mid X, Z) = \sum_{i=1}^n \tilde{Z}_{ik} [x_i\log \lambda_k - \lambda_k] + \text{constant}.
\]

Differentiating with respect to $\lambda_k$ gives
\[
\frac{\partial \ell_c}{\partial \lambda_k}
=
\sum_{i=1}^n \tilde{Z}_{ik}
\left[
\frac{x_i}{\lambda_k}-1
\right].
\]

At the optimum,
\[
\frac{\partial \ell_c}{\partial \lambda_k}=0.
\]

Therefore,
\[
\sum_{i=1}^n \tilde{Z}_{ik}
\left[
\frac{x_i}{\lambda_k}-1
\right]
=0.
\]

Expanding the sum gives
\[
\frac{1}{\lambda_k}\sum_{i=1}^n \tilde{Z}_{ik}x_i
-
\sum_{i=1}^n \tilde{Z}_{ik}
=0.
\]

Therefore,
\[
\frac{1}{\lambda_k}\sum_{i=1}^n \tilde{Z}_{ik}x_i
=
\sum_{i=1}^n \tilde{Z}_{ik}.
\]

Multiplying both sides by $\lambda_k$ gives
\[
\sum_{i=1}^n \tilde{Z}_{ik}x_i
=
\lambda_k
\sum_{i=1}^n \tilde{Z}_{ik}.
\]

Therefore,
\[
\lambda_k
=
\frac{\sum_{i=1}^n \tilde{Z}_{ik}x_i}
{\sum_{i=1}^n \tilde{Z}_{ik}}.
\]

Since
\[
n_k=\sum_{i=1}^n\tilde{Z}_{ik},
\]
we can write
\[
\boxed{
\hat{\lambda}_k
=
\frac{1}{n_k}
\sum_{i=1}^n\tilde{Z}_{ik}x_i
}.
\]

Thus, the estimated Poisson parameter for group $k$ is the
mean of the observations belonging to group $k$.

\paragraph{Step 2: Verify that the stationary point is a maximum.}

The second derivative with respect to $\lambda_k$ is
\[
\frac{\partial^2\ell_c}{\partial\lambda_k^2}
=
-\sum_{i=1}^n
\tilde{Z}_{ik}
\frac{x_i}{\lambda_k^2}.
\]

Since
\[
\tilde{Z}_{ik}\geq 0,
\qquad
x_i\geq 0,
\qquad
\lambda_k>0,
\]
we have
\[
\frac{\partial^2\ell_c}{\partial\lambda_k^2}
\leq 0.
\]

Therefore, the stationary point corresponds to a maximum.

Hence,
\[
\boxed{
\hat{\lambda}_k
=
\frac{\sum_{i=1}^n\tilde{Z}_{ik}x_i}
{\sum_{i=1}^n\tilde{Z}_{ik}}
}.
\]

That is, the MLE of $\lambda_k$ is the sample mean of the
observations assigned to component $k$.
